\section{Boolean gluing of extensions}
\label{glue:section}

We prove that the configuration in Proposition~\ref{ext:reduction}
cannot exist.  Throughout this section, $T$ is the target simple group
of the homomorphisms, whereas $R$ is the simple type of the domain
socles.  We retain the notation $A,D,H,U_J,\beta_J,C_J,N_J,I_J$ of
that proposition.  All extension domains lie in the finite group $D$.

The extensions will glue if their domain cores have images containing
$\operatorname{Inn}(T)$.  We first prove this criterion: a least normal
subgroup of $A$ with that image is normalized by every domain, and
conjugation on its simple quotient supplies the common extension.
We then force the criterion on the Boolean faces.  For an atom $x$,
the normal image $\alpha(C_x)$ is either trivial or contains the inner
group.  In the trivial case, compatible lifts from the three rank-two
faces through $x$ produce a different common extension, again forbidden
by \ref{ext:e4}.  Hence every atom core retains the inner group;
intersections pass this property to the larger proper faces.

\subsection{A criterion for extending over generated domains}

We begin with the elementary fact that makes least subgroups with a
prescribed perfect image available.

\begin{lemma}[Intersections retaining a perfect image]
\label{glue:perfect-intersections}
Let $f:M\to H_0$ be a homomorphism of finite groups, let $P\le H_0$
be perfect, and let $B_1,\ldots,B_s\trianglelefteq M$, with $s\ge1$, satisfy
$P\le f(B_i)$ for every $i$.  Then
\[
                  P\le f\left(\bigcap_{i=1}^s B_i\right).
\]
\end{lemma}
\begin{proof}
For two normal subgroups,
\[
 f(B_1\cap B_2)\ \supseteq\ f([B_1,B_2])
   =[f(B_1),f(B_2)]\ \supseteq\ [P,P]=P.
\]
Induction proves the assertion for the finite family.
\end{proof}

The following criterion applies to any finite family of extension
domains.  It will be used at the end with two rank-three Boolean faces
whose union is full.

\begin{lemma}[A common core and its conjugation extension]
\label{glue:common-core}
Let $A\le U_i\le D$ be a finite nonempty family of subgroups of a finite
group.  Let $T$ be nonabelian simple and suppose that homomorphisms
\[
 \alpha:A\to\operatorname{Aut}(T),\qquad
 \beta_i:U_i\to\operatorname{Aut}(T)
\]
satisfy
\[
 \beta_i|_A=\alpha,
 \qquad
 \alpha\bigl(\operatorname{core}_{U_i}(A)\bigr)
                         \ge\operatorname{Inn}(T).
\]
Put $K=\langle U_i:i\rangle$.  Then
$\alpha(\operatorname{core}_K(A))\ge\operatorname{Inn}(T)$, and there
is a unique homomorphism $\beta:K\to\operatorname{Aut}(T)$ extending
all the $\beta_i$.
\end{lemma}
\begin{proof}
Write $J=\operatorname{Inn}(T)$ and let
$\iota:T\to J$ be the isomorphism $t\mapsto c_t$.
Consider
\[
 Q=\bigcap\{B\trianglelefteq A:\alpha(B)\ge J\}.
\]
The family is nonempty.  Since $J$ is perfect,
Lemma~\ref{glue:perfect-intersections} shows that its finite
intersection still maps over $J$.  The subgroup $\alpha^{-1}(J)$
belongs to the family, so in fact
\begin{equation}\label{glue:least-perfect-image}
                         \alpha(Q)=J.
\end{equation}
Thus $Q$ is the least normal subgroup of $A$ with image containing $J$.

Fix $i$ and put $C_i=\operatorname{core}_{U_i}(A)$.  The hypothesis
gives $Q\le C_i$.  Make the analogous construction in $C_i$:
\[
 Q_i=\bigcap\{B\trianglelefteq C_i:\alpha(B)\ge J\}.
\]
The same intersection argument, using $C_i\cap\alpha^{-1}(J)$,
gives $\alpha(Q_i)=J$.
For $u\in U_i$ and $c\in C_i$,
\[
       \alpha(ucu^{-1})
          =\beta_i(u)\alpha(c)\beta_i(u)^{-1}.
\]
Since $C_i\trianglelefteq U_i$ and
$J\trianglelefteq\operatorname{Aut}(T)$, conjugation by $U_i$ permutes
the defining family for $Q_i$.  Thus $Q_i\trianglelefteq U_i$, and in
particular $Q_i\trianglelefteq A$.  Minimality of $Q$ gives $Q\le Q_i$.
Conversely, $Q\le C_i$, $Q\trianglelefteq C_i$, and
\eqref{glue:least-perfect-image} give $Q_i\le Q$.
Consequently $Q_i=Q$ for every $i$.

It follows that $Q\trianglelefteq K$ and
$Q\le\operatorname{core}_K(A)$, proving the common-core assertion.
Set $L=Q\cap\ker\alpha$.  Since
$L=Q\cap\ker\beta_i$ for every $i$, the subgroup $L$ is normalized by
each $U_i$ and therefore by $K$.
The map
\[
 \vartheta:Q/L\longrightarrow T,
        \qquad \vartheta(qL)=\iota^{-1}(\alpha(q))
\]
is an isomorphism.  Transporting the conjugation action of $K$ on
$Q/L$ through $\vartheta$ defines a homomorphism
$\beta:K\to\operatorname{Aut}(T)$.
For $u\in U_i$ and $q\in Q$,
\[
 \alpha(uqu^{-1})
 =\beta_i(u)\alpha(q)\beta_i(u)^{-1}
 =\iota\bigl(\beta_i(u)(\vartheta(qL))\bigr).
\]
As $\vartheta$ is onto, this proves $\beta|_{U_i}=\beta_i$.
These restrictions determine $\beta$ uniquely because the $U_i$
generate $K$.
\end{proof}

\subsection{An atom whose core is killed}

For $B\subseteq I_J$, let
\begin{equation}\label{glue:support-preimage}
 N_J(B)=q_J^{-1}\left(\prod_{i\in B}S_{J,i}\right),
 \qquad
 q_J:N_J\longrightarrow S_J=N_J/C_J,
\end{equation}
where $S_{J,i}$ are the simple direct factors of $S_J$.
Thus $N_J(\varnothing)=C_J$.  For comparable nonempty index sets
$J\subset K$, write $p_{KJ}:I_K\to I_J$ for the canonical factor map.

\begin{lemma}\label{glue:factorization}
If $x$ is an original Boolean atom and $\alpha(C_x)=1$, then
$\beta_x$ factors through the permutation action of $U_x$ on $I_x$.
\end{lemma}
\begin{proof}
By \ref{ext:e3}, $U_x=A\ker\beta_x$.  Since $A<U_x$ and $C_x\le A$,
the kernel is not contained in $C_x$.  The assumption makes
$\beta_x$ a homomorphism on $U_x/C_x$, so its nontrivial kernel there
contains the unique minimal normal subgroup $N_x/C_x$.
Consequently $\beta_x(N_x)=1$.

Let $P$ be the kernel of the action of $U_x$ on $I_x$.
Modulo $C_x$, the faithful conjugation action on $S_x$ embeds
$P/N_x$ into $\operatorname{Out}(R)^{I_x}$.  To justify faithfulness,
the centralizer of $S_x$ in $U_x/C_x$ is normal and disjoint from
$S_x$; uniqueness of the nonabelian minimal normal subgroup makes it
trivial.  Schreier's Theorem therefore makes $P/N_x$ soluble.
Its image under $\beta_x$ is a soluble normal subgroup of
$H=\alpha(A)$.  Lemma~\ref{ext:almost-simple} makes that image trivial,
so $P\le\ker\beta_x$, as claimed.
\end{proof}

Fix an atom $x$.  To rule out $\alpha(C_x)=1$, we shall show that this assumption
produces a subgroup $L$ normalized by $A$, generated by subgroups on
which the local extensions are trivial.  The remaining
condition for the formula $\theta(a\ell)=\alpha(a)$ to define an
extension on $AL$ is $A\cap L\le\ker\alpha$.  By Lemma~\ref{glue:factorization}, it is
enough to prove that every element of $A\cap L$ fixes the factor set
$I_x$.  The three proper faces through $x$ provide the local generators;
their proper pairwise joins will let us compare different factor labels.

Set
\[
 \mathcal U_x=\{x\cup\{a\}:a\in\{1,\ldots,4\}\setminus x\}.
\]
This family has three members.  Distinct $u,v\in\mathcal U_x$ have a
proper join $w=u\cup v$ of rank three.  For $j\in I_x$ and
$y\supseteq x$, write
\[
                   F_{y,j}=p_{yx}^{-1}(j),
       \qquad F_{x,j}=\{j\}.
\]
All factor maps in this notation are defined on proper nodes.
Figure~\ref{fig:boolean-faces} displays these three faces and their
proper pairwise joins.

A lift of the fibre $F_{u,j}$ must lie over the corresponding support
$F_{u\cup v,j}$ in each proper pair join $u\cup v$.  We therefore intersect its support
preimages in all those joins.  The perfect-image intersection lemma
ensures that these simultaneous requirements do not lose any factor
of the desired fibre.

\begin{lemma}[Compatible lifts of the fibres]
\label{glue:lifts}
For $u\in\mathcal U_x$ and $j\in I_x$, define
\begin{equation}\label{glue:Luj}
 L_{u,j}=N_u(F_{u,j})\cap
       \bigcap_{\substack{v\in\mathcal U_x\\v\ne u}}
       \bigl(U_u\cap N_{u\cup v}(F_{u\cup v,j})\bigr).
\end{equation}
These subgroups have the following properties.
\begin{enumerate}[label=\textup{(\roman*)}]
\item They are normal in $N_u$, and
$q_u(L_{u,j})=\prod_{i\in F_{u,j}}S_{u,i}$.
\item Conjugation by $U_x$ permutes the choices according to its action
on $I_x$.  In particular, $C_x$ normalizes every $L_{u,j}$.
\item $U_u=\langle A,L_{u,j}:j\in I_x\rangle$.
\item For $j\ne k$, one has
\[
 [L_{u,j},L_{v,k}]\le
 \begin{cases}
 C_{u\cup v},&u\ne v,\\
 C_u,&u=v.
 \end{cases}
\]
Both right-hand sides are contained in $C_x$.
\item If $\alpha(C_x)=1$, then $\beta_u(L_{u,j})=1$ for all $u,j$.
\end{enumerate}
\end{lemma}
\begin{proof}
Fix a proper pair join $w=u\cup v$.  Equal-type transport gives
\begin{equation}\label{glue:Nu-transport}
                  N_u=(U_u\cap N_w)C_u.
\end{equation}
The first factor on the right acts by inner automorphisms on
$N_w/C_w$, hence preserves every coordinate subproduct there.
The second factor satisfies $C_u\le C_x\le U_x$.  Since $C_x$ acts
trivially on $I_x$, equivariance of $p_{wx}$ says that $C_u$ fixes the
set $F_{w,j}$.  It too normalizes $N_w(F_{w,j})$.
Thus $U_u\cap N_w(F_{w,j})$ is normal in $N_u$ and lies in $N_u$.

In $N_w/C_w$, the subdirect group $(U_u\cap N_w)/C_w$ has its diagonal
strips indexed by $I_u$.  Their supports are the fibres of $p_{wu}$.
Because $p_{wx}=p_{ux}p_{wu}$, intersecting with the $F_{w,j}$-support
product selects exactly the strips indexed by $F_{u,j}$.
The isomorphism from this subdirect group to $N_u/C_u$ therefore shows
that $U_u\cap N_w(F_{w,j})$ maps onto
$\prod_{i\in F_{u,j}}S_{u,i}$ under $q_u$.
The first group in \eqref{glue:Luj} is normal in $N_u$ and has the same
image.  This image is perfect.  Lemma~\ref{glue:perfect-intersections}
proves (i).

Every domain occurring in \eqref{glue:Luj} contains $U_x$, and all the
factor maps are $U_x$-equivariant.  Its conjugation therefore permutes
the whole definition with $j$, proving (ii).  The images in (i)
generate $N_u/C_u$, and $C_u\le A$; since $U_u=AN_u$, this proves (iii).

For $u\ne v$, the two lifts in (iv) lie in disjoint coordinate support
products of $N_{u\cup v}/C_{u\cup v}$.  They commute in that quotient.
For $u=v$ the identical argument takes place in $N_u/C_u$.
Core inclusion along $U_x\le U_u\le U_{u\cup v}$ gives the asserted
containment in $C_x$.

Finally suppose that $\alpha(C_x)=1$.  Then $\beta_u(C_u)=1$.
The kernel of $\beta_u$ supplements $A$ and is not contained in $C_u$,
so it contains $N_u$ by the unique-minimal-normal-subgroup argument in
Lemma~\ref{glue:factorization}.  This proves (v).
\end{proof}

We next construct subgroups $B_k$ which distinguish the labels
$k\in I_x$ and are normalized by every lift with a different label.
Different labels commute only modulo $C_x$, and $C_x$ need not be
normal in the group generated by all the lifts.  The following finite
core argument keeps a full simple image without requiring that
normality.

\begin{lemma}[Taking cores without losing a simple image]
\label{glue:retained-core}
Let $R_0\le D$, let $C\trianglelefteq R_0$, and suppose that
$R_0/C$ is nonabelian simple.  Let $L_1,\ldots,L_s\le D$ be normalized
by $R_0$.  Suppose that, for every $t$, there is a subgroup
$D_t\trianglelefteq R_0$ such that
\[
 D_tC=R_0,\qquad D_t^{L_t}=D_t,
 \qquad [D_t,L_t]\le D_t\cap C.
\]
Then, for $W=\langle R_0,L_1,\ldots,L_s\rangle$, the subgroup
$\operatorname{core}_W(R_0)$ maps onto $R_0/C$.
No normalization of $C$ by $L_t$ is assumed.
\end{lemma}
\begin{proof}
Let $\pi:R_0\to R_0/C$ be the quotient map.  Start with $B=R_0$.
At every stage $B$ will be normal in $R_0$ and satisfy
$\pi(B)=R_0/C$.  For a fixed $t$, replace $B$ by
\begin{equation}\label{glue:core-operation}
       F_t(B)=\bigcap_{\ell\in L_t}(B\cap D_t)^\ell,
       \qquad K^\ell=\ell^{-1}K\ell.
\end{equation}
The subgroup $B\cap D_t$ still maps onto $R_0/C$ by
Lemma~\ref{glue:perfect-intersections}.  Put $K=B\cap D_t$.  The subgroup $K$ is normal in $D_t$, because
$B\trianglelefteq R_0$ and $D_t\leq R_0$.  Since $L_t$ normalizes
$D_t$, every conjugate $K^\ell$ lies in $D_t$ and is normal in $D_t$.
For $d\in D_t$ and $\ell\in L_t$, the assumption
$[D_t,L_t]\leq D_t\cap C$ gives $\pi(d^\ell)=\pi(d)$.
Consequently $\pi(K^\ell)=\pi(K)=R_0/C$ for every $\ell$.
This applies the quotient map only to elements of $D_t\leq R_0$;
it does not require $C$ to be normalized by $L_t$.  Applying
Lemma~\ref{glue:perfect-intersections} inside $D_t$ shows that $F_t(B)$
has full image.

The new subgroup is normal in $R_0$.  Indeed, $R_0$ normalizes $B$, $D_t$, and $L_t$.  For $r\in R_0$
one therefore has $(K^\ell)^r=(K^r)^{\ell^r}=K^{\ell^r}$, and
$\ell^r\in L_t$.  Thus $R_0$ permutes precisely the conjugates
intersected in \eqref{glue:core-operation}.  The operation is descending, since the
identity conjugate occurs.  Apply $F_1,\ldots,F_s$ in repeated rounds.
Finiteness forces stabilization, and at a stable round every
intermediate inclusion is equality.  The resulting subgroup $B_*$
is therefore normalized by each $L_t$ as well as by $R_0$.  Thus
$B_*\trianglelefteq W$ and $B_*\le R_0$, while
$\pi(B_*)=R_0/C$.  Since
$B_*\le\operatorname{core}_W(R_0)$, the asserted core also has full
image.
\end{proof}

\begin{lemma}[Subgroups detecting the factor labels]
\label{glue:detectors}
Put $L_j=\langle L_{u,j}:u\in\mathcal U_x\rangle$ and
$R_k=N_x(\{k\})$.  There are subgroups $B_k\le R_k$, for $k\in I_x$,
such that
\begin{enumerate}[label=\textup{(\roman*)}]
\item $q_x(B_k)=S_{x,k}$;
\item every $L_j$ with $j\ne k$ normalizes $B_k$;
\item $C_x$ normalizes every $B_k$, and $A$ permutes the $B_k$ in the
same way as it permutes $I_x$.
\end{enumerate}
In particular, the subgroups $B_k$ are pairwise distinct.
\end{lemma}
\begin{proof}
For $u\in\mathcal U_x$ set
\[
                     D_{u,k}=U_x\cap N_u(F_{u,k}).
\]
The transported subdirect group $(U_x\cap N_u)/C_u$ maps isomorphically
onto $N_x/C_x$.  Selecting its strips over $k$ shows that
\begin{equation}\label{glue:Duk}
 D_{u,k}\le R_k,\quad D_{u,k}C_x=R_k,
 \quad D_{u,k}\cap C_x=C_u.
\end{equation}
It is normal in $R_k$: all elements of $R_k\le N_x$ fix the factor
labels $I_x$, and the support-preimage construction is equivariant.

For $j\ne k$, both $D_{u,k}$ and $L_{u,j}$ lie in $N_u$ and have
disjoint supports modulo $C_u$.  Thus
\[
                    [D_{u,k},L_{u,j}]\le C_u.
\]
Since $C_u\le D_{u,k}$, this also proves that $L_{u,j}$ normalizes
$D_{u,k}$.  Moreover $R_k$ normalizes $L_{u,j}$ by
Lemma~\ref{glue:lifts}(ii), since $R_k$ fixes every label of $I_x$.
We can therefore apply Lemma~\ref{glue:retained-core} to $R_k$, its
normal subgroup $C_x$, and all the lifts $L_{u,j}$ with $j\ne k$, using
the subgroup $D_{u,k}$ for the lift labelled by $(u,j)$.

Define canonically
\[
 W_k=\langle R_k,L_j:j\ne k\rangle,
 \qquad B_k=\operatorname{core}_{W_k}(R_k).
\]
That lemma proves (i); the definition proves (ii).
Since $C_x\le R_k\le W_k$, it normalizes $B_k$.
Conjugation by $A$ permutes the pairs $(W_k,R_k)$ and hence their
cores, proving (iii).  Their images in $N_x/C_x$ are the distinct
nontrivial coordinate factors, so the $B_k$ are distinct.
\end{proof}

To use the subgroups $B_k$, we must collect all letters with one
chosen label at the end of a word in the lifts.  The errors lie in
$C_x$, which normalizes each individual lift.  The following word
calculation records exactly the form needed.

\begin{lemma}[Collection with labelled factors]
\label{glue:collection}
Let $C\le D$ normalize subgroups $P_\lambda\le D$ indexed by a finite
set.  Assign each index a label $c(\lambda)\in I$.  Suppose that
\[
 [P_\lambda,P_\mu]\le C
                \quad\text{whenever }c(\lambda)\ne c(\mu).
\]
Put $L_j=\langle P_\lambda:c(\lambda)=j\rangle$ and
$L=\langle L_j:j\in I\rangle$.  For every $k\in I$, every element
$g\in L$ can be written
\[
        g=pqc,\qquad
        p\in\langle L_j:j\ne k\rangle,\quad q\in L_k,\quad c\in C.
\]
\end{lemma}
\begin{proof}
Write $g$ as a word whose letters belong to the individual
$P_\lambda$; inverses are allowed because these are subgroups.
Interchange adjacent letters with different labels to move the
$k$-labelled letters to the end.  Each interchange introduces a
commutator in $C$.  Move that error to the right through the remaining
word, using
\[
                    c\ell=(c\ell c^{-1})c.
\]
The conjugated letter remains in its original $P_\lambda$, so its
label does not change.  Each interchange reduces the number of pairs
consisting of a $k$-letter preceding a non-$k$-letter.  The process
terminates with the asserted expression.  This is a word calculation;
it does not form a quotient by $C$ or require $C\trianglelefteq L$.
\end{proof}

\begin{proposition}\label{glue:atom-cores}
For every original Boolean atom $x$,
$\alpha(C_x)\ge\operatorname{Inn}(T)$.
\end{proposition}
\begin{proof}
The subgroup $\alpha(C_x)$ is normal in
$\beta_x(U_x)=H$.  Lemma~\ref{ext:almost-simple} reduces the assertion
to excluding $\alpha(C_x)=1$.  Assume that equality and use the lifts
and detecting subgroups just constructed.  Put
\[
                         L=\langle L_j:j\in I_x\rangle.
\]
The group $A$ normalizes $L$.

Let $a\in A\cap L$.  We show that $a$ fixes every member of $I_x$.
For subgroup conjugation use the right-action notation
$K^g=g^{-1}Kg$.  If $a$ sent $k$ to $h\ne k$, equivariance would give
$B_k^a=B_h$.  Lemmas~\ref{glue:lifts} and~\ref{glue:collection}, with
$C=C_x$ and the individual $L_{u,j}$ as labelled subgroups, give
\[
        a=pqc,\qquad
        p\in\langle L_j:j\ne k\rangle,\quad q\in L_k,\quad c\in C_x.
\]
The element $p$ normalizes $B_k$, $q$ normalizes $B_h$, and $c$
normalizes both.  It follows successively that
\[
                 B_k^{qc}=B_h,
       \qquad B_k^q=B_h,
       \qquad B_k=B_h,
\]
contradicting Lemma~\ref{glue:detectors}.  Thus $a$ fixes the factor
set.  Lemma~\ref{glue:factorization} now gives $\alpha(a)=1$.

Consequently $A\cap L\le\ker\alpha$.  Because $A$ normalizes $L$,
the formula
\[
              \theta:AL\longrightarrow\operatorname{Aut}(T),
                       \qquad \theta(a\ell)=\alpha(a)
\]
defines a homomorphism.  It is well defined by the intersection
inclusion; multiplication is respected since $L\trianglelefteq AL$.
For each $u\in\mathcal U_x$, Lemma~\ref{glue:lifts}(iii),(v) says that
$U_u$ is generated by $A$ and lifts killed by both $\theta$ and
$\beta_u$.  Hence $\theta|_{U_u}=\beta_u$.

The three sets $u\in\mathcal U_x$ have full union
$\{1,\ldots,4\}$.  Their generated domain is $AL$, so this common
extension contradicts \ref{ext:e4}.  The assumption
$\alpha(C_x)=1$ is impossible.
\end{proof}

\subsection{Compatible joins and the contradiction}

\begin{lemma}[Cores under a compatible join]
\label{glue:core-joins}
For every nonempty proper original Boolean node $z$,
\[
             C_z=\bigcap_{x\subseteq z,\,|x|=1}C_x.
\]
Consequently $\alpha(C_z)\ge\operatorname{Inn}(T)$.
\end{lemma}
\begin{proof}
The assertion is immediate for an atom, so suppose that $|z|\ge2$.
Every $C_x$ in the intersection contains $C_z$.  Pass to $U_z/C_z$
and use bars for the images of $A,U_x,C_x$.  Put
\[
 N=N_z/C_z,\qquad E_x=(U_x\cap N_z)/C_z.
\]
The transport isomorphism in \ref{ext:e2} gives
$E_x\cap\overline C_x=1$.  Both are normal in $\overline U_x$, so
$[E_x,\overline C_x]=1$.

Since $U_z=AN_z$ and $A\le U_x$, one has
$\overline U_x=\overline A E_x$.
Let $M=\langle E_x:x\subseteq z,\ |x|=1\rangle$.
Each $E_x$ is normalized by $\overline A$.  The compatible-join
generation in \ref{ext:e2} therefore gives
\[
                       \overline U_z=\overline A M,
\]
so $M\trianglelefteq\overline U_z$.  It is a nontrivial subgroup of
the minimal normal subgroup $N$, and hence $M=N$.
Thus $\bigcap_x\overline C_x$ centralizes $N$.
The centralizer is normal in $\overline U_z$ and intersects $N$ in
$Z(N)=1$.  If it were nontrivial it would contain the unique minimal
normal subgroup $N$, a contradiction.  The intersection is therefore
trivial, which proves the equality of cores.

Each $C_x$ is normal in $A$ and maps over $\operatorname{Inn}(T)$ by
Proposition~\ref{glue:atom-cores}.  Apply
Lemma~\ref{glue:perfect-intersections} with
$P=\operatorname{Inn}(T)$ to obtain the image assertion.
\end{proof}

\begin{theorem}\label{glue:contradiction}
The Boolean extension configuration of
Proposition~\ref{ext:reduction} does not exist.
\end{theorem}
\begin{proof}
Take the two proper Boolean nodes
\[
              z=\{1,2,3\},\qquad w=\{2,3,4\}.
\]
Lemma~\ref{glue:core-joins} gives
$\alpha(C_z),\alpha(C_w)\ge\operatorname{Inn}(T)$.
Lemma~\ref{glue:common-core} therefore extends $\beta_z$ and $\beta_w$
to $\langle U_z,U_w\rangle$.  Since $z\cup w=\{1,\ldots,4\}$, this
contradicts \ref{ext:e4}.
\end{proof}

\begin{proof}[Proof of Theorem~\ref{main:group}]
Choose the finite decorated lattice of
Proposition~\ref{lat:decorated-lattice}.  If either it or its dual had a
finite-group interval representation, Proposition~\ref{ext:reduction}
would produce the configuration excluded by
Theorem~\ref{glue:contradiction}.  Hence this lattice has no such
representation, proving Theorem~\ref{main:group}.
\end{proof}

Theorem~\ref{main:congruences} now follows from
Theorem~\ref{main:group} and the global P\'alfy--Pudl\'ak equivalence
\cite[Theorem~2]{PalfyPudlak1980}.
